% mazzi: 42,43
\chapter{La programmazione della preparazione fisica nel calcio}

% ---------------------------------------------------------------------------
\section{Gli interrogativi del preparatore fisico}

La domanda di fondo è \textbf{come progettare l'uso degli esercizi} nel ciclo annuale, nei
macrocicli, nei microcicli e nella singola seduta. Da lì si diramano gli interrogativi concreti:

\begin{itemize}
  \item che \textbf{tipo di esercizi} scegliere: generali, speciali o specifici?
  \item \textbf{quando} proporli?
  \item \textbf{quanta} e \textbf{quale} espressione di forza serve alla prestazione?
  \item quanto e come rispettare l'\textbf{accordo dinamico} con l'azione tecnica (o
    tecnico-tattica)?
  \item come \textbf{combinare} o \textbf{mettere in sequenza} gli esercizi di forza con quelli di
    carattere tecnico-tattico --- il problema della cosiddetta ``trasformazione'', \textbf{che non
    esiste};
  \item come essere sicuri che l'incremento della performance \textbf{non produca danni} invece che
    vantaggi (\textbf{effetto boomerang}), e non solo in termini di infortuni ma anche di
    interferenze negative sulla tecnica specifica.
\end{itemize}

\rubrica{Perché la ``trasformazione'' non esiste}
In natura non si trasforma nulla: parlare di esercitazioni di forza e poi, sul campo, di
``esercizi di trasformazione'' è una dicitura scorretta. La formulazione corretta sarebbe
\textbf{esercizi di riadattamento neuromuscolare}.

% ---------------------------------------------------------------------------
\section{L'allenamento non è una scienza esatta}

\subsection{I tempi cambiano}

Tre confronti d'epoca introducono il tema:

\begin{itemize}
  \item \textbf{i giovani}: nel 1970 i genitori dovevano tirare i figli per le orecchie per farli
    \textit{rientrare} a casa; nel 2017 la situazione è capovolta e devono tirarli per le orecchie
    per farli \textit{uscire}, staccandoli dal divano e dai videogiochi;
  \item \textbf{lo studio}: prima esclusivamente su libri e in biblioteca, sul materiale cartaceo;
    oggi su internet;
  \item \textbf{la formazione del tecnico}: con YouTube, Facebook, Instagram, Pinterest e Twitter
    allenatori e preparatori rischiano di diventare dei \textbf{\textit{social sport trainer}}.
\end{itemize}

\rubrica{Ma i principi dell'allenamento no}
\begin{itemize}
  \item \textbf{principi di base}: dosabilità, variabilità, trasferibilità;
  \item \textbf{effetti desiderati}:
    \begin{itemize}
      \item performance $\rightarrow$ effetti \textbf{prestativi};
      \item pre-infortunio $\rightarrow$ effetti \textbf{preventivi};
      \item post-infortunio $\rightarrow$ effetti \textbf{riabilitativi};
    \end{itemize}
  \item \textbf{background fisiologico};
  \item \textbf{evidenza scientifica}.
\end{itemize}

Fisiologia e biomeccanica sono sempre le stesse: non cambiano.

\subsection{Un ambito contraddittorio (Nicola Maffiuletti)}

Dal \textit{Human Performance Lab} della Schulthess Clinic di Zurigo:

\begin{itemize}
  \item ambito \textbf{contraddittorio e mutevole};
  \item \textbf{poche evidenze}, e quindi infinite teorie e possibilità --- soprattutto negli sport
    di squadra;
  \item \textbf{miti, mode, imitazioni, ciarlatani};
  \item il \textbf{buon senso}, che si è perso.
\end{itemize}

\subsection{Allenamento in ``laboratorio'' ed evidenze scientifiche}

Negli ultimi \textbf{20--25 anni}: boom della ricerca applicata, pertinenza e utilità spesso dubbie,
``scientificità'' da chiarire. La \textbf{piramide dell'evidenza}, dal livello più alto al più
basso:

\begin{enumerate}
  \item recensioni sistematiche;
  \item studi randomizzati e controllati;
  \item studi longitudinali: studio di un gruppo di persone osservando gli eventi che si verificano
    \textbf{senza alcun intervento} da parte dello sperimentatore;
  \item studi di controllo su casi;
  \item serie di casi;
  \item casi clinici;
  \item opinioni di esperti.
\end{enumerate}

\subsection{L'allenamento ``funzionale''}

L'allenamento è \textbf{sempre stato funzionale}, o meglio deve esserlo: qualsiasi programma deve
mettere l'atleta nelle condizioni ideali per esprimere la propria performance. Negli ultimi anni,
però, sulla spinta delle \textbf{ditte specializzate}, l'etichetta ha assunto un significato
ristretto:

\begin{itemize}
  \item esercizi \textbf{pluri-articolari} e \textbf{multiplanari}, accelerazioni e decelerazioni,
    stabilità, \textit{core};
  \item impiego in \textbf{preparazione} e in \textbf{stagione};
  \item obiettivo dichiarato: miglioramento della \textbf{qualità del movimento} a fine prestativo e
    preventivo.
\end{itemize}

\rubrica{La presa di coscienza}
L'allenamento funzionale è \textbf{difficilmente controllabile} nel dosaggio e
\textbf{moderatamente efficace}. Tre attrezzi lo mostrano bene:

\begin{itemize}
  \item l'\textbf{elastico}: offre la sua massima resistenza \textbf{al termine della corsa}, quindi
    non è un mezzo valido per allenare forza esplosiva, reattività e rapidità;
  \item il \textbf{TRX}: c'è stato un periodo in cui senza di esso non si poteva essere preparatori
    fisici;
  \item il \textbf{gesto tecnico su superficie instabile}: diverse ricerche indicano che lavorare
    così non allena la prevenzione né la tecnica, ma \textbf{peggiora la qualità esecutiva} del
    gesto.
\end{itemize}

\subsection{Alcuni consigli}

\begin{itemize}
  \item $\downarrow$ materiali;
  \item $\uparrow$ buon senso;
  \item $\uparrow$ attenzione sui giovani e sulla loro costruzione;
  \item $\uparrow$ qualità scientifica di ciò che si progetta;
  \item approccio \textbf{mirato ma semplice};
  \item $\uparrow$ fattori \textbf{neuro-cognitivi};
  \item $\uparrow$ \textbf{esplosività} e \textbf{tecnica fine}.
\end{itemize}

% ---------------------------------------------------------------------------
\section{L'atleta al centro del processo allenante}

Le aree che concorrono alla prestazione sono nove: \textbf{allenamento}, \textbf{valutazione},
\textbf{gare}, \textbf{nutrizione}, \textbf{medicina dello sport}, \textbf{psicologia dello sport},
\textbf{recupero}, \textbf{riabilitazione}, \textbf{prevenzione}, con la \textbf{ricerca} a fare da
raccordo.

Lo stesso schema si presta a due letture opposte, e il confronto fra le due è il punto:

\begin{itemize}
  \item la deriva: tutte queste aree finiscono \textbf{subissate dai dati numerici} che arrivano dai
    preparatori, e sono i dati a occupare il centro;
  \item l'impostazione corretta: al centro sta l'\textbf{atleta}, e i professionisti responsabili
    delle varie aree \textbf{interagiscono fra loro} perché l'atleta abbia il massimo supporto
    nell'esprimere la propria performance.
\end{itemize}

% ---------------------------------------------------------------------------
\section{Allenamento e allenabilità}

\subsection{Un po' di storia}

Da \textit{La preparazione fisica ottimale del calciatore} di J.~Weineck --- libro datato ma sempre
attuale:

\begin{itemize}
  \item ``Il segreto del successo nel calcio sta sempre nell'allenamento'' (Beenhakker, 1990);
  \item ``Il miglior maestro per l'allenamento è la gara'' (Cramer, 1987);
  \item ``Dalla gara capiamo che cosa dobbiamo allenare'' (Krauspe/Rauhut/Teschner, 1990).
\end{itemize}

\subsection{Che cos'è l'allenamento sportivo}

\textbf{Allenamento sportivo}: processo d'azione complesso che ha lo scopo di influire, in modo
pianificato e rivolto a un obiettivo specifico, sul livello di prestazione sportiva e sulla capacità
di realizzare nel migliore dei modi possibile tale prestazione.

\begin{itemize}
  \item \textbf{processo d'azione complesso}: indirizzato a raggiungere effetti adeguati sugli
    \textbf{indici rilevanti} della prestazione dell'atleta;
  \item \textbf{pianificato}: obiettivi, metodi, contenuti, costruzione e organizzazione
    dell'allenamento vengono \textbf{prestabiliti}, tenendo conto sia delle nozioni della scienza
    dell'allenamento sia dell'esperienza.
\end{itemize}

\subsection{L'allenabilità e le condizioni della prestazione}

\textbf{Allenabilità}: grado di adattamento ai carichi di allenamento. È un parametro che dipende da
una serie di fattori endogeni ed esogeni.

\rubrica{Condizioni della prestazione dirette endogene}
\begin{itemize}
  \item \textbf{condizione}: capacità di forza, capacità di resistenza, capacità di rapidità,
    mobilità articolare;
  \item \textbf{tecnica dei movimenti}: capacità coordinative, abilità motorie;
  \item \textbf{tattica}: capacità di risolvere problemi in modo adeguato alla situazione.
\end{itemize}

\rubrica{Condizioni della prestazione indirette endogene}
\begin{itemize}
  \item \textbf{costituzione}: sistema muscolare, sistema scheletrico e legamentoso;
  \item \textbf{capacità organiche-fisiche}: sistema respiratorio, sistema cardiocircolatorio;
  \item \textbf{capacità cognitive}: sistema ormonale;
  \item \textbf{capacità emotive}: fegato, reni e organi digerenti, sistema nervoso e organi di
    senso.
\end{itemize}

% ---------------------------------------------------------------------------
\section{Le dimensioni e l'organizzazione del processo allenante}

\subsection{Le dimensioni}

Tutto discende dallo \textbf{stile di allenamento} (\textit{training style}) che si sceglie, che si
articola in due dimensioni:

\begin{itemize}
  \item \textbf{gestione della squadra} (\textit{team management});
  \item \textbf{metodologia di lavoro} (\textit{working methodology}), che a sua volta si divide in:
    \begin{itemize}
      \item \textbf{allenamento collettivo} $\rightarrow$ la \textbf{squadra};
      \item \textbf{allenamento individuale} $\rightarrow$ il \textbf{singolo calciatore}.
    \end{itemize}
\end{itemize}

\subsection{Lo schema sintetico del processo allenante}

Lo schema d'insieme parte da due domande --- \textit{che cosa voglio?} e \textit{come lo ottengo?}
--- che si traducono nel \textbf{modello di gioco} e nello \textbf{stile di allenamento}, in dialogo
fra loro. Da lì si scende alla \textbf{metodologia di allenamento} e alla \textbf{gestione della
squadra}, e quindi ai due binari dell'allenamento collettivo e individuale. La
\textbf{periodizzazione} li organizza su due livelli:

\begin{itemize}
  \item \textbf{micro-struttura}: la \textbf{seduta} --- articolata in pre-sessione, parte iniziale,
    parte principale (classificata per complessità e dinamica degli sforzi, con i relativi compiti),
    parte finale e post-sessione --- dentro il \textbf{microciclo} settimanale;
  \item \textbf{macro-struttura}: i \textbf{cicli} o \textbf{blocchi} dentro la
    \textbf{stagione}, da luglio a maggio.
\end{itemize}

\subsection{L'organizzazione della seduta}

\rubrica{La struttura del lavoro giornaliero}
\begin{enumerate}
  \item \textbf{allenamento individuale}: pre-sessione e pre-attivazione, su progettazione
    individuale;
  \item \textbf{allenamento collettivo}: la seduta sul campo;
  \item \textbf{allenamento individuale}: post-sessione e recupero.
\end{enumerate}

\rubrica{La struttura della seduta collettiva}
\begin{enumerate}
  \item \textbf{parte iniziale}: riscaldamento generale, riscaldamento specifico, attivazione
    neuromuscolare;
  \item \textbf{parte principale}: compiti classificati per \textbf{complessità} e per
    \textbf{dinamica degli sforzi};
  \item \textbf{parte finale}: \textit{cool down}, collegata alle attività post-allenamento.
\end{enumerate}

\subsection{I tipi di riscaldamento}

\begin{itemize}
  \item \textbf{senza palla}: onde, quadrati, linee e figure, giochi;
  \item \textbf{con palla}: piccoli gruppi, quadrati, sequenze di passaggi, giochi, esercitazioni
    tattiche.
\end{itemize}

\subsection{La complessità dell'allenamento}

I compiti della parte principale si classificano su \textbf{quattro livelli di complessità}
crescente:

\begin{enumerate}
  \item \textbf{condizionale}: esercitazioni \textbf{senza palla} e con presa di decisione
    \textbf{non specifica};
  \item \textbf{tecnico}: \textbf{con palla}, coordinazione motoria specifica e presa di decisione
    non specifica;
  \item \textbf{tattico}: il \textbf{processo decisionale} del calciatore viene stimolato. Si
    articola a sua volta in quattro gradini:
    \begin{enumerate}
      \item \textbf{esercizi di passaggio}: si muove prioritariamente la \textbf{palla}, non i
        giocatori;
      \item \textbf{esercizi di possesso}: si muovono \textbf{giocatori e palla}, senza gioco di
        posizione;
      \item \textbf{esercitazioni posizionali}: giocatori e palla in movimento, con
        \textbf{gioco di posizione} e \textbf{direzione di gioco} (1--2 fasi di gioco);
      \item \textbf{giochi posizionali}: giocatori e palla in movimento, gioco di posizione e giochi
        a direzione con \textbf{tutte le fasi} e con vincoli;
    \end{enumerate}
  \item \textbf{competitivo}: \textbf{pochi vincoli}, applicazione dei principi modellati sulla
    competizione.
\end{enumerate}

\subsection{La dinamica degli sforzi}

Alle tre \textbf{zone di intensità} corrispondono cinque \textbf{dinamiche dello sforzo}, che sono
la scala con cui si scompone il volume di ogni categoria di esercitazione:

\begin{table}[htbp]
  \centering
  \small
  \renewcommand{\arraystretch}{1.3}
  \begin{tabularx}{\textwidth}{|l|l|X|}
    \hline
    \textbf{Intensità} & \textbf{Dinamica} & \textbf{Descrizione} \\
    \hline
    Bassa & Bassa intensiva & Attività condotte molto \textbf{sotto} l'intensità di competizione ---
      recupero attivo, spesso usate come \textit{cool down} \\
    \hline
    Media & Media estensiva & Attività condotte sotto l'intensità di competizione, \textbf{non
      rigenerative} \\
    \hline
    \multirow{3}{*}{Alta} & Lunga intensiva & Attività alla stessa intensità della competizione o
      \textbf{superiore}, e di \textbf{lunga durata} \\
    \cline{2-3}
     & Corta intensiva & Attività \textbf{sopra} l'intensità di competizione, di \textbf{breve
      durata} e con \textbf{recupero incompleto} \\
    \cline{2-3}
     & Massima intensità & Attività alla \textbf{massima intensità}, breve durata e
      \textbf{recupero maggiore} \\
    \hline
  \end{tabularx}
\end{table}

Le cinque dinamiche si distinguono anche graficamente, mettendo l'intensità in ordinata e il volume
in ascissa: al crescere della dinamica i picchi si alzano (da \textit{low} verso \textit{maximal}) e
si diradano, cioè il numero delle ripetizioni si riduce mentre ciascuna diventa più intensa.

\rubrica{Quattro esempi di allenamento fisico}
\begin{itemize}
  \item \textbf{bassa intensità} --- $24'$: corsetta sul lato lungo del campo, camminata per
    mezzo campo, di nuovo corsetta a fondo campo, camminata, corsetta, e così via;
  \item \textbf{media intensità}: circuito di forza e coordinazione, con corse più dinamiche e
    piccoli cambi di direzione;
  \item \textbf{media intensità e lunga durata} --- $24'$: circuiti condizionali mirati ad agilità e
    coordinazione --- rapidità dei piedi sui cerchi, allungo a superare il paletto a centrocampo,
    ostacoli, allungo, ostacoli saltati a piedi pari, allungo con la scaletta;
  \item \textbf{alta intensità e corta durata}: esplosività, velocità e coordinazione degli arti
    inferiori, su un circuito di \textbf{quattro stazioni} ---
    \begin{enumerate}
      \item \textit{skip} su scaletta e sprint;
      \item \textit{skip} di agilità sui cerchi e sprint laterale, senza cambi di direzione;
      \item cambi di direzione sui paletti;
      \item salti sugli ostacoli a una gamba con cambi di direzione repentini, ostacoli e uscita in
        sprint.
    \end{enumerate}
\end{itemize}

% ---------------------------------------------------------------------------
\section{La programmazione nei professionisti: il peso dei dati}

Nei professionisti la programmazione non può prescindere dai dati che arrivano dalle strumentazioni:
\textit{training load}, GPS, RPE, Borg, CPK, VAS, TQR, profili individuali, previsione e
\textit{scouting} --- il mondo dei \textbf{big data} applicati alla squadra.

\subsection{Il volume di allenamento per categoria di esercitazione}

Il volume, misurato in \textbf{minuti}, viene ripartito fra le categorie di esercitazione lungo i
microcicli successivi. Le categorie sono:

\begin{itemize}
  \item \textbf{riscaldamento senza palla};
  \item \textbf{riscaldamento con palla};
  \item \textbf{condizionale};
  \item \textbf{tecnico};
  \item \textbf{tattico};
  \item \textbf{competitivo}.
\end{itemize}

Ciascuna delle ultime quattro è a sua volta scomposta per \textbf{intensità}: bassa intensità,
estensiva, lunga-intensiva, breve-intensiva, massima intensità.

\subsection{Il confronto con le altre squadre}

I dati di squadra, confrontati con quelli di tutto il campionato, si leggono su due piani: la
distanza totale percorsa nel primo e nel secondo tempo, e la sua ripartizione per fasce di velocità
--- cammino ($0$--$2$\,m/s, cioè $0$--$7{,}2$\,km/h), corsa leggera ($2$--$4$\,m/s,
$7{,}2$--$14{,}4$\,km/h), corsa quantitativa ($4$--$5{,}5$\,m/s, $14{,}4$--$19{,}8$\,km/h),
accelerazioni ($5{,}5$--$7$\,m/s, $19{,}8$--$25{,}2$\,km/h) e sprint ($> 7$\,m/s,
$> 25{,}2$\,km/h).

Il valore diagnostico sta nello \textbf{scarto dalla media del campionato}. In un caso reale la
squadra risultava sopra la media in quasi tutti i parametri, ma \textbf{sotto} --- di poco --- nella
distanza percorsa in \textbf{sprint} oltre i $25{,}2$\,km/h: un segnale preciso che qualcosa nel
processo di allenamento andava modificato.

\subsection{L'indice di usura (\textit{stress index})}

Indicatore semplice, adottabile da chiunque:

\[
\text{Stress Index} = \frac{\text{minuti giocati}}{\text{minuti delle partite}} \times 100
\]

\begin{itemize}
  \item oltre il \textbf{75\%}: il calciatore è sottoposto a un \textit{training load} e a un
    \textit{overuse} molto elevati, per l'eccesso di minutaggio;
  \item sotto il \textbf{50\%}: il problema si ribalta --- sono i giocatori che stanno
    costantemente in panchina o in tribuna, che vanno tenuti in condizione e resi pronti alla
    richiesta della gara. Per loro, oltre all'aspetto fisico, conta quello \textbf{psicologico}: la
    scarsa disponibilità al surplus di allenamento che devono sopportare.
\end{itemize}

% ---------------------------------------------------------------------------
\section{La metodologia dell'allenamento}

\begin{itemize}
  \item la qualità del processo allenante può essere ottimizzata se si seguono i principi della
    \textbf{specificità} e dell'\textbf{individualità};
  \item sia le \textbf{capacità fisiche specifiche} sia le \textbf{abilità tecniche} possono essere
    migliorate impiegando mezzi allenanti \textbf{sport-specifici};
  \item un'\textbf{affinità coordinativa} fra esercizi di allenamento e gara ha il vantaggio di
    portare lo stimolo di adattamento proprio nei distretti muscolari e articolari che determinano
    il gesto specifico.
\end{itemize}

\subsection{Sport specific training}

\begin{itemize}
  \item si riferisce a esercizi ad \textbf{alta intensità di movimento} che \textbf{simulano il
    gesto tecnico} utilizzato in un determinato tipo di attività sportiva;
  \item uno studio ha dimostrato che allenamenti di forza con \textbf{esercizi imitativi} del gesto
    tecnico migliorano questo tipo di abilità (Wilson e coll., 1996).
\end{itemize}

\rubrica{Il presupposto dell'alta intensità}
L'alta intensità di movimento richiede che l'atleta abbia già una \textbf{buona qualità tecnica}: se
nell'esercitazione si commettono errori bisogna fermarsi e abbassare l'intensità, e a quel punto
questa tipologia di allenamento \textbf{perde la sua valenza}.

\subsection{Esercizi generali, speciali e specifici}

È la distinzione che risponde a uno degli interrogativi di partenza. I tre tipi si definiscono per
la \textbf{distanza dal gesto di gara} e hanno ciascuno un obiettivo proprio.

\begin{table}[htbp]
  \centering
  \small
  \renewcommand{\arraystretch}{1.3}
  \begin{tabularx}{\textwidth}{|l|X|X|}
    \hline
    \textbf{Tipo} & \textbf{Che cosa sono} & \textbf{A che cosa servono} \\
    \hline
    \textbf{Generali} & Esercizi che non presentano \textbf{alcun elemento} del gesto tecnico della
      disciplina e che si discostano per tempo di esecuzione, posizione e spostamento rispetto al
      gesto di gara & Sviluppo della \textbf{forza massima ed esplosiva} a carattere generale \\
    \hline
    \textbf{Speciali} & Esercizi che \textbf{rispettano il gesto di gara} ma modificano le
      caratteristiche \textbf{spazio-temporali} della tecnica, riducendo o aumentando la velocità
      rispetto al gesto di gara & Migliorare la coordinazione \textbf{intra}- e
      \textbf{intermuscolare} per perfezionare la tecnica \\
    \hline
    \textbf{Specifici} & Esercizi che \textbf{corrispondono} agli esercizi di gara, in condizioni
      \textbf{vicine o uguali} alla competizione & \textbf{Stabilizzare la tecnica} attraverso la
      ripetizione sistematica dei gesti \\
    \hline
  \end{tabularx}
\end{table}

\rubrica{Esercizi generali}
\begin{itemize}
  \item \textbf{senza sovraccarico}: piegamenti a terra, dorsali, salto della funicella, addominali;
  \item \textbf{con sovraccarico}: il lavoro alle macchine --- \textit{lat machine},
    \textit{bench press}, \textit{calf machine}, \textit{leg press}, \textit{abductor} e
    \textit{adductor machine}, \textit{leg curl}, \textit{hip rotary} --- e il lavoro con i
    \textbf{bilancieri liberi}, da preferire.
\end{itemize}

\rubrica{Esercizi speciali}
\begin{itemize}
  \item \textbf{senza sovraccarico}: sprint, sprint in salita e in discesa, le varie tipologie di
    \textit{skip} e di corse calciate;
  \item \textbf{con sovraccarico}: traino, balzi con sovraccarico, pliometria nelle varie forme e
    anche in modo \textbf{asimmetrico}.
\end{itemize}

\rubrica{Esercizi specifici}
\begin{itemize}
  \item tecnica individuale;
  \item percorsi nelle varie tipologie, a 1, a 2, a 3, e situazioni $1 > 1$, $2 > 2$ e simili;
  \item tutte le esercitazioni tecnico-tattiche di competenza del mister e degli staff tecnici;
  \item gare amichevoli e ufficiali.
\end{itemize}

\rubrica{Un caso di confine}
Partire seduti su una panca e staccare subito di testa per colpire la palla, o scendere da un plinto
per andare a colpirla: esercitazioni di questo tipo sono \textbf{generali}, non speciali né
specifiche, perché non hanno nulla a che vedere con la specificità del gesto tecnico.

% ---------------------------------------------------------------------------
\section{Il carico di allenamento}

\subsection{Le definizioni}

\begin{itemize}
  \item \textbf{carico di allenamento}: somma del lavoro richiesto all'atleta, ovvero l'insieme
    delle \textbf{sollecitazioni funzionali};
  \item la stessa nozione dal lato degli effetti: sommatoria di tutti gli effetti sull'organismo
    dell'atleta che provocano il \textbf{cambiamento del suo stato funzionale}
    (Y.~V.~Verchoshanskij, 2001).
\end{itemize}

\rubrica{Potenziale ed effetto allenante}
\begin{itemize}
  \item \textbf{potenziale allenante del carico}: la \textbf{possibilità} di produrre una reazione
    funzionale di adattamento dell'organismo dell'atleta e i corrispondenti cambiamenti del suo
    stato funzionale;
  \item \textbf{effetto allenante del carico} (EA): il \textbf{risultato} dell'influenza esercitata
    dal carico sull'organismo, che si esprime nel cambiamento \textbf{concreto} dello stato
    funzionale dell'atleta.
\end{itemize}

Il primo è potenzialità, il secondo è ciò che si verifica davvero.

\subsection{Carico esterno e carico interno}

\begin{itemize}
  \item \textbf{carico esterno}: è tutto quello che si può \textbf{misurare} --- p.\,es.\ chilometri
    percorsi, velocità di percorrenza, tempi di recupero;
  \item \textbf{carico interno}: è la \textbf{somma degli stress} a cui viene sottoposto l'organismo
    dell'atleta.
\end{itemize}

\subsection{La specificità del carico allenante}

\begin{itemize}
  \item \textbf{intensità}: concetto che esprime il \textbf{livello} (la percentuale) dell'impegno
    richiesto al soggetto rispetto alle sue \textbf{capacità massimali};
  \item \textbf{volume}: insieme delle caratteristiche numeriche, della durata e del numero delle
    \textbf{ripetizioni} che lo stimolo assume in una unità di allenamento;
  \item \textbf{densità}: rapporto fra il \textbf{tempo di lavoro} e il \textbf{tempo di recupero}
    dell'unità --- recupero fra un'esercitazione e l'altra --- o del ciclo di allenamento, p.\,es.\
    quello settimanale.
\end{itemize}

\subsection{Carico di allenamento: volume $\times$ intensità}

\rubrica{Il volume}
Si valuta attraverso le \textbf{unità di misura} delle attività proposte all'atleta; è compito degli
staff tecnici programmarne il peso.

\rubrica{Le fasce di carico (da Viru, modificato)}
In percentuale sul \textbf{massimo carico sopportabile}:

\begin{itemize}
  \item fino al \textbf{20\%}: carico \textbf{insufficiente};
  \item \textbf{20--40\%}: \textbf{recupero};
  \item \textbf{40--60\%}: \textbf{mantenimento};
  \item \textbf{60--80\%}: \textbf{sviluppo} prestativo;
  \item \textbf{80--100\%}: carico \textbf{eccessivo}, che può produrre più problemi che vantaggi.
\end{itemize}

\rubrica{L'intensità: due modi di classificarla}
\begin{enumerate}
  \item \textbf{scala generale}, quella più spesso usata dai tecnici: leggera, media, forte,
    intensa, massimale. Il limite evidente del metodo è la \textbf{soggettività};
  \item \textbf{percentuale rispetto a un parametro conosciuto}. Prendendo per esempio
    l'allenamento di resistenza:
    \begin{itemize}
      \item la frequenza cardiaca massima ($FC_{max}$);
      \item la velocità aerobica massimale (VAM);
      \item la velocità di soglia anaerobica ($V_{obla}$);
      \item la velocità record sulla distanza proposta.
    \end{itemize}
\end{enumerate}

Sta al preparatore fisico decidere da quale metodo partire per quantificare e valutare l'intensità
di ciò che propone.

% ---------------------------------------------------------------------------
\section{Pianificare, periodizzare, programmare}

Sono i tre atti dell'organizzazione del processo allenante, in ordine di scala decrescente:

\begin{enumerate}
  \item \textbf{pianificazione}: formulazione della \textbf{strategia}, attraverso tappe successive,
    di differenti tipi di carico in un \textbf{ampio spazio di tempo} (annuale, pluriennale);
  \item \textbf{periodizzazione}: formulazione di \textbf{principi teorici} relativi a
    \textbf{periodi più particolareggiati} dell'intera pianificazione;
  \item \textbf{programmazione}: applicazione dei principi teorici della periodizzazione, cioè la
    \textbf{stesura del programma} di allenamento.
\end{enumerate}

Nel calcio la pianificazione vera e propria spesso \textbf{non avviene}: si resta legati al
risultato più che alla progettazione a lungo termine.

% ---------------------------------------------------------------------------
\section{La periodizzazione}

\subsection{Le tappe di una struttura periodizzata}

Adattato da Issurin (2010):

\begin{table}[htbp]
  \centering
  \small
  \renewcommand{\arraystretch}{1.3}
  \begin{tabularx}{\textwidth}{|l|l|X|}
    \hline
    \textbf{Tappa} & \textbf{Durata} & \textbf{Contenuto} \\
    \hline
    Preparazione pluriennale & Anni & Piano sistematico annuale o pluriennale, sviluppato su cicli di
      2 o 4 anni \\
    \hline
    Macrociclo & Mesi & Grande ciclo di allenamento; include i periodi di preparazione, competizione
      e transizione \\
    \hline
    Mesociclo & Settimane & Ciclo di media ampiezza, costituito da un certo numero di microcicli \\
    \hline
    Microciclo & Giorni & Piccolo ciclo di allenamento costituito da un certo numero di giorni,
      frequentemente una settimana \\
    \hline
    Seduta di allenamento & Minuti/ore & Singola seduta, svolta individualmente o in gruppo \\
    \hline
  \end{tabularx}
\end{table}

\subsection{La periodizzazione adattata al calcio}

\begin{itemize}
  \item \textbf{macrociclo annuale}: i 12 mesi dell'anno;
  \item \textbf{periodo agonistico}: dalla metà di \textbf{luglio} a tutto il mese di \textbf{maggio}
    dell'anno successivo --- il confine varia con i campionati che la squadra disputa e con il
    calendario delle federazioni;
  \item \textbf{periodo transitorio}: quello che resta, nei professionisti di norma \textbf{un mese o
    un mese e mezzo};
  \item \textbf{mesocicli}: i mesi, da luglio a luglio;
  \item \textbf{microcicli}: le settimane in cui ogni mese viene suddiviso per la stesura del
    programma.
\end{itemize}

\subsection{I tre periodi e la loro caratteristica}

\begin{itemize}
  \item \textbf{periodo preparatorio}: sviluppo \textbf{ininterrotto} della condizione fisica,
    dall'inizio alla fine del periodo --- ininterrotto perché ogni giorno, e dove è possibile con
    doppia seduta giornaliera, si produce qualcosa mirato al miglioramento della condizione;
  \item \textbf{periodo competitivo}: sviluppo della condizione fisica \textbf{in forma sportiva} ---
    quella che serve a rispondere alle richieste fisiologiche della gara --- e sua
    \textbf{stabilizzazione} nel corso del periodo;
  \item \textbf{periodo di transizione}: la \textbf{non prevedibilità} del livello di scadimento
    della condizione fisica.
\end{itemize}

% ---------------------------------------------------------------------------
\section{La progressione metodologica}

L'ordine con cui i contenuti entrano nella preparazione si è \textbf{capovolto}:

\begin{itemize}
  \item l'impostazione di molti anni fa partiva dalle \textbf{corse di lunghissima durata} ad
    andatura quasi costante;
  \item l'impostazione attuale procede \textbf{tecnica $\rightarrow$ forza $\rightarrow$
    resistenza}: si comincia subito dalle esercitazioni \textbf{tecniche specifiche}, si passa alla
    \textbf{forza specifica}, e la \textbf{resistenza} diventa una conseguenza di quanto fatto nelle
    fasi precedenti, non un obiettivo diretto di partenza.
\end{itemize}

È la stessa \textbf{inversione della piramide} già descritta a proposito della storia della
preparazione fisica (§\ref{sec:inversione-piramide}), qui letta sull'asse
\textbf{fibre lente} (base) $\rightarrow$ \textbf{fibre rapide} (vertice): sono le fibre rapide a
contare di più per la performance del calciatore, e si interviene su quelle per prime.

% ---------------------------------------------------------------------------
\section{I tempi di recupero}

Tempi necessari per realizzare il recupero a seguito dei vari tipi di allenamento:

\begin{table}[htbp]
  \centering
  \small
  \renewcommand{\arraystretch}{1.3}
  \begin{tabularx}{\textwidth}{|X|c|c|c|c|c|}
    \hline
    \textbf{Tipo di allenamento} & \textbf{Res.\ aerobica} & \textbf{Res.\ anaerobica} &
      \textbf{Forza esplosiva} & \textbf{Ipertrofia} & \textbf{Velocità tecnica} \\
    \hline
    Recupero incompleto & --- & 1,5--2 h & 2--3 h & 2--3 h & 2--3 h \\
    \hline
    Recupero quasi completo & 12 h & 12 h & 12--18 h & 18 h & 18 h \\
    \hline
    Recupero completo & 24--36 h & 24--48 h & 48--72 h & 72--84 h & 72 h \\
    \hline
  \end{tabularx}
\end{table}

Il valore operativo della tabella sta nel decidere \textbf{a distanza di quanto} si può riproporre un
contenuto:

\begin{itemize}
  \item la \textbf{resistenza aerobica} --- il lavoro delle prime fasi di preparazione estiva --- ha
    recupero quasi completo a \textbf{12 ore}, quindi intermittente e ripetute si possono riproporre
    anche il \textbf{giorno successivo};
  \item la \textbf{resistenza anaerobica} e la \textbf{forza esplosiva} non si possono riproporre se
    non dopo \textbf{due o tre giorni} dalla seduta.
\end{itemize}

% ---------------------------------------------------------------------------
\section{La preparazione precampionato}

\subsection{Che cos'è}

\begin{itemize}
  \item \textbf{base fondamentale} dell'intera stagione agonistica;
  \item ha caratteristiche di \textbf{unicità, originalità e inimitabilità};
  \item \textbf{presupposto essenziale}: l'approfondita conoscenza degli atleti;
  \item \textbf{obiettivo}: il conseguimento di un elevato grado di adattamenti a stimoli
    \textbf{crescenti}, di tipo sia fisico sia tecnico-tattico;
  \item necessita di \textbf{programmazione e coordinamento}: non deve essere un periodo
    circoscritto e isolato, ma un progetto coordinato su tutto l'arco della stagione;
  \item non deve essere un mero \textbf{``contenitore'' di esercitazioni}.
\end{itemize}

\subsection{Gli obiettivi generali}

\begin{itemize}
  \item \textbf{psicologico}: conoscere i nuovi giocatori, il loro carattere e la loro personalità;
    amalgamare i nuovi acquisti col resto della squadra; creare un gruppo che interpreti lo spirito
    di squadra;
  \item \textbf{tecnico-tattico}: migliorare la tecnica individuale dei singoli; sviluppare e
    definire un modulo di gioco con l'attuazione degli schemi;
  \item \textbf{fisico}:
    \begin{itemize}
      \item migliorare le \textbf{qualità aerobiche} dei singoli, sia a livello centrale sia
        periferico;
      \item migliorare il livello di \textbf{forza}, in particolare degli arti inferiori, senza
        trascurare la muscolatura posturale e gli arti superiori;
      \item aumentare la \textbf{resistenza alla velocità} e \textbf{alla rapidità};
      \item ottimizzare la \textbf{prevenzione} di tipo propriocettivo.
    \end{itemize}
\end{itemize}

\subsection{L'andamento di volume e intensità}

Sull'arco delle \textbf{6--7 settimane} che portano al periodo competitivo, volume e intensità si
scambiano il peso:

\begin{itemize}
  \item \textbf{settimane 1--2}, \textit{general fitness}: predomina il \textbf{volume}, a discapito
    dell'intensità;
  \item \textbf{settimane 3--6}, sviluppo della \textbf{fitness specifica}: intorno alla quinta o
    sesta settimana avviene l'\textbf{inversione} --- il volume scema e l'intensità delle
    esercitazioni aumenta;
  \item \textbf{settimana 7}, \textit{preparation}: l'intensità è al suo massimo, il volume al
    minimo, e si entra nella settimana tipo del periodo competitivo.
\end{itemize}

% ---------------------------------------------------------------------------
\section{Il microciclo in periodo agonistico}

I contenuti sono \textbf{indicazioni di massima}, non prescrizioni.

\subsection{Cinque allenamenti settimanali (gara al sabato)}

\begin{table}[htbp]
  \centering
  \small
  \renewcommand{\arraystretch}{1.3}
  \begin{tabularx}{\textwidth}{|l|l|X|}
    \hline
    \textbf{Giorno} & \textbf{Seduta} & \textbf{Contenuto} \\
    \hline
    Lunedì & 1\textsuperscript{a} & Potenza anaerobica e alcune esercitazioni di forza \\
    \hline
    Martedì & 2\textsuperscript{a} & Massimo carico di forza \\
    \hline
    Mercoledì & 3\textsuperscript{a} & Massima qualità tecnica e situazioni \\
    \hline
    Giovedì & 4\textsuperscript{a} & Velocità, rapidità e forza elastica \\
    \hline
    Venerdì & 5\textsuperscript{a} & Palle inattive e richiamo di rapidità \\
    \hline
    Sabato & --- & \textbf{Gara} \\
    \hline
    Domenica & --- & Riposo \\
    \hline
  \end{tabularx}
\end{table}

Totali: \textbf{5} allenamenti settimanali, \textbf{20 circa} mensili, \textbf{200 circa} annuali su
circa 10 mesi.

\subsection{Quattro allenamenti settimanali (gara alla domenica)}

\begin{table}[htbp]
  \centering
  \small
  \renewcommand{\arraystretch}{1.3}
  \begin{tabularx}{\textwidth}{|l|l|X|}
    \hline
    \textbf{Giorno} & \textbf{Seduta} & \textbf{Contenuto} \\
    \hline
    Lunedì & --- & Riposo \\
    \hline
    Martedì & 1\textsuperscript{a} & Potenza anaerobica e alcune esercitazioni di forza \\
    \hline
    Mercoledì & 2\textsuperscript{a} & Massima qualità tecnica e situazioni \\
    \hline
    Giovedì & 3\textsuperscript{a} & Velocità e rapidità \\
    \hline
    Venerdì & 4\textsuperscript{a} & Forza elastica e palle inattive \\
    \hline
    Sabato & --- & Riposo \\
    \hline
    Domenica & --- & \textbf{Gara} \\
    \hline
  \end{tabularx}
\end{table}

Totali: \textbf{4} allenamenti settimanali, \textbf{16 circa} mensili, \textbf{160 circa} annuali.

\subsection{Il progetto annuale}

Riportare il microciclo su un calendario di dodici mesi permette di contare gli allenamenti
effettivi. Le convenzioni del calendario:

\begin{itemize}
  \item \textbf{giugno} e \textbf{luglio}: periodo transitorio;
  \item \textbf{agosto}: periodo di preparazione;
  \item \textbf{giornate in giallo}: i cicli settimanali del periodo competitivo;
  \item \textbf{caselle con la data in azzurro}: le feste tradizionali, quindi le pause natalizie.
\end{itemize}

Il conteggio dà \textbf{circa 183} allenamenti annuali nell'ipotesi a cinque sedute e
\textbf{circa 145} in quella a quattro, a cui vanno \textbf{aggiunte le partite}: amichevoli, di
campionato e di coppa.

% ---------------------------------------------------------------------------
\section{La settimana di allenamento: due modelli a confronto}

\subsection{Programmazione ``italiana'' e ``scozzese''}

\begin{table}[htbp]
  \centering
  \small
  \renewcommand{\arraystretch}{1.3}
  \begin{tabularx}{\textwidth}{|l|X|X|}
    \hline
    \textbf{Giorno} & \textbf{``Italiana'' (gara domenica)} & \textbf{``Scozzese'' (gara sabato)} \\
    \hline
    Lunedì & Riposo & 90$'$ \\
    \hline
    Martedì & 75$'$ --- allenamento di ripresa & 100$'$ \\
    \hline
    Mercoledì & 160$'$ --- quasi sempre doppia seduta & Riposo \\
    \hline
    Giovedì & 60$'$ --- partitella tecnico-tattica con la squadra B, le giovanili o una squadra di
      categoria inferiore & 100$'$ \\
    \hline
    Venerdì & 80$'$ --- rapidità e reattività & 75$'$ \\
    \hline
    Sabato & 60$'$ --- tattica di gioco & 100$'$ + \textbf{gara} \\
    \hline
    Domenica & 100$'$ + \textbf{gara} & Riposo \\
    \hline
    \textbf{Totale} & \textbf{535$'$} & \textbf{465$'$} \\
    \hline
  \end{tabularx}
\end{table}

\rubrica{Che cosa cambia}
\begin{itemize}
  \item il modello scozzese ha un \textbf{minutaggio settimanale complessivo minore} ($465'$ contro
    $535'$) ma \textbf{due giorni di riposo} --- domenica e mercoledì;
  \item nonostante ciò il \textbf{minutaggio quotidiano è molto più elevato} e distribuito in modo
    uniforme ($90'$--$100'$ al giorno), mentre l'italiano concentra il carico in
    un'unica grande seduta di metà settimana;
  \item soprattutto, in Scozia \textbf{non esiste l'allenamento di ripristino} per chi ha giocato:
    al primo allenamento della settimana chi ha disputato la gara si allena alla \textbf{stessa
    intensità} di chi non l'ha giocata, altrimenti il \textit{manager} non lo prende in
    considerazione per la partita successiva. La \textbf{motivazione indotta dall'allenatore} è
    quindi essa stessa un dato della programmazione.
\end{itemize}

\subsection{Quanto spazio ha davvero la preparazione fisica}

Riepilogo settimanale di una stagione di Serie B (2011/12):

\begin{itemize}
  \item \textbf{preparazione fisica}: 115 minuti, il \textbf{15,54\%};
  \item \textbf{allenamento tecnico-tattico}: 625 minuti, l'\textbf{84,46\%};
  \item di cui alla \textbf{forza} solo il \textbf{6,75\%}, cioè circa \textbf{30 minuti a
    settimana}.
\end{itemize}

È il dato che ridimensiona ogni discorso sull'allenamento della forza \textbf{a secco}: il vincolo
non è metodologico ma di \textbf{tempo disponibile}.

% ---------------------------------------------------------------------------
\section{L'analisi di una sessione allenante}

\subsection{Il quadro globale}

Seduta reale, monitorata con cardiofrequenzimetro:

\begin{itemize}
  \item \textbf{durata totale}: 95 minuti;
  \item $FC_{max}$ \textbf{173} bpm, $FC$ media \textbf{133} bpm, $FC_{min}$ \textbf{81} bpm;
  \item spesa energetica \textbf{468,6} kcal, indice di lavoro \textbf{202,1};
  \item tempo per zona: $Z_1$ $27'36''$, $Z_2$ $10'20''$, $Z_3$ $7'40''$, $Z_4$ $25'36''$,
    $Z_5$ $12'52''$.
\end{itemize}

La distribuzione ha la sua moda fra i \textbf{105 e i 123 bpm}: letta così, la seduta
\textbf{non è ottimale} rispetto alle richieste fisiologiche della partita. È però il dato
aggregato a ingannare --- la seduta era modulata in \textbf{tre blocchi} molto diversi fra loro.

\subsection{L'analisi per tipologia di carico}

\begin{table}[htbp]
  \centering
  \small
  \renewcommand{\arraystretch}{1.3}
  \begin{tabularx}{\textwidth}{|X|c|c|c|}
    \hline
    \textbf{Blocco} & \textbf{Tempo a FC bassa} & \textbf{Tempo a FC alta} & \textbf{Durata} \\
    \hline
    Torelli & $16'20''$ & $1'52''$ & $\approx 18'$ \\
    \hline
    Partite a cavallo di metà campo & $5'16''$ & $26'12''$ & $\approx 31'$ \\
    \hline
    $300$\,m $\times$ 4 ripetute $\times$ 2 serie & $15'54''$ & $10'24''$ & $\approx 26'$ \\
    \hline
  \end{tabularx}
\end{table}

\begin{itemize}
  \item \textbf{torelli}: intensità molto bassa, con l'atleta a lungo in \textbf{zona 1} (fra i 70 e
    i 105 bpm);
  \item \textbf{partite a cavallo di metà campo}: l'atleta resta molto tempo in \textbf{zona 4} e in
    parte anche in \textbf{zona 5}, vicino alla frequenza cardiaca massima;
  \item \textbf{ripetute sui 300 m}: il recupero obbligato fra una ripetizione e l'altra produce un
    \textbf{decadimento} della frequenza cardiaca --- $15'54''$ dei $26'$ totali stanno nelle zone 1,
    2 e 3 e solo $10'24''$ sono ripetizioni effettive. Ne risulta una scarsa \textbf{specificità}
    rispetto al lavoro calcistico.
\end{itemize}

\subsection{La conclusione}

Confrontando la distribuzione per zone dei tre blocchi, la \textbf{zona 4} è quella che separa
netto: le \textbf{partite} la occupano in modo massiccio, torelli e ripetute quasi per nulla. La
\textbf{partita} è quindi il mezzo più idoneo e fondamentale per il miglioramento delle qualità
fisiche del calciatore.
